\begin{table} \begin{small} \begin{center}\begin{threeparttable}[b] \caption{Migration rates by origin, Indonesia} \label{Migration rates by origin, Indonesia} \begin{tabular}{l*{3}c} \toprule
                    &       Rural&       Urban&         All\\
\midrule
\multicolumn{4}{l}{\emph{1995}} \\Migration rate      &        32.3&        35.8&        33.7\\
Moves within category&        31.1&        74.6&        49.4\\
\multicolumn{4}{l}{\emph{2011}} \\Migration rate      &        38.7&        33.7&        35.7\\
Moves within category&        24.4&        84.2&        58.7\\
\multicolumn{4}{l}{\emph{2012}} \\Migration rate      &        38.9&        34.1&        35.8\\
Moves within category&        25.4&        83.8&        60.7\\
\bottomrule \end{tabular} \label{} \begin{tablenotes} \item {\footnotesize \emph{Notes: } Data source:  1995 Supas; 2011 Susenas; 2012 Susenas. Migration is measured as living in a regency other than the birth regency. Regencies are classified as rural or urban based on the share of their population that report being rural; we choose the cutoff to classify the regency as rural to match the national urbanization rate for each year. } \end{tablenotes} \end{threeparttable} \end{center} \end{small} \end{table}
